\documentclass[11pt]{amsart}
\usepackage{amssymb}
\usepackage[hidelinks]{hyperref}
\input{glyphtounicode}
\pdfgentounicode=1
\newtheorem{theorem}{Theorem}[section]
\newtheorem{lemma}[theorem]{Lemma}
\newtheorem{proposition}[theorem]{Proposition}
\newtheorem{fact}{Fact}
\renewcommand{\thefact}{F\arabic{fact}}
\numberwithin{equation}{section}
\DeclareMathOperator{\rad}{rad}
\newcommand{\Z}{\mathbb Z}
\newcommand{\Q}{\mathbb Q}
\newcommand{\R}{\mathbb R}
\newcommand{\F}{\mathbb F}
\newcommand{\eps}{\varepsilon}

\newcommand{\fixedparameters}{%
\begin{table}[ht]
\caption{Parameters fixed by \(e=\min(\eps,1/10)\).}\label{tab:parameters}
\centering
\renewcommand{\arraystretch}{1.25}
\begin{tabular}{ll}
\hline
Arithmetic scales & \(g=e/100,\quad e_0=e^2/10,\quad c_0=e/10\)\\
Degree and width & \(k=\lceil2/e\rceil,\quad w=e/(k+1)\)\\
Dimensions & \(m=\lceil2(k+1)/w\rceil,\quad d=m+1\)\\
Fourier cutoff & \(\eta=w/[4(k+1)],\quad K=\lceil8m/(\pi^2\eta^{m+1})\rceil\)\\
Degree scale & \(\kappa=e_0/(d+3)\)\\
\hline
\end{tabular}
\end{table}}

\title[A subexponential upper bound for log c]{A subexponential upper bound for log c in the abc conjecture}
\author{Ruiyi Zhang}
\address{Independent Researcher, Columbia, Maryland, USA}
\email{ryzfsu@gmail.com}
\date{November 20, 2026}
\subjclass[2020]{Primary 11D75; Secondary 11J86, 11L20}
\keywords{abc conjecture, radical, interpolation determinant, exponential sum over primes}
\hypersetup{
  pdftitle={A subexponential upper bound for log c in the abc conjecture},
  pdfauthor={Ruiyi Zhang},
  pdfsubject={2020 Mathematics Subject Classification: Primary 11D75, Secondary 11J86 and 11L20},
  pdfkeywords={abc conjecture, radical, interpolation determinant, exponential sum over primes}
}
\begin{document}
\begin{abstract}
Let \(a,b,c\) be positive integers with \(a+b=c\) and \(\gcd(a,b)=1\).
We prove the bound
\(\log c\le C_\varepsilon\rad(abc)^\varepsilon\) for every
\(\varepsilon>0\). Under the contrary height assumption, a prime power
dividing one side forces a nonzero polynomial to vanish on part of an
integer simplex. A second prime splits that set into residue classes.
Finite differences and Lucas's congruence show that its points determine
every polynomial up to the required degree.
\end{abstract}
\maketitle
\section{Introduction}

Write \(R=\rad(abc)=\prod_{p\mid abc}p\) and \(H=\log c\).
The \(abc\) conjecture predicts \(H\le(1+\varepsilon)\log R+O_\varepsilon(1)\)
for every coprime positive solution of \(a+b=c\). It remains open.
Stewart and Yu obtained \(H\ll R^{1/3}(\log R)^3\)
\cite[Theorem~1]{SY2001}. Pasten and Sep\'ulveda-Manzo obtained
subexponential estimates under size conditions on the smaller summand
\cite[Theorems~1.2 and~1.3]{PastenSep}. Zhang gives
\(H\ll R^{1/3}\exp(-\log R/(3000\log\log R))\) for all sufficiently large
\(R\) \cite[Theorem~1]{ZhangThird} and subsequently obtains
\(H\ll_\varepsilon R^{2/9+\varepsilon}\) \cite[Theorem~1.1]{Z29}.

\begin{theorem}\label{thm:main}
Let \(\varepsilon>0\). There is a constant \(C_\varepsilon>0\) such that
\[
       \log c\le C_\varepsilon\rad(abc)^\varepsilon
\]
for all coprime positive integers \(a,b\) and \(c=a+b\).
\end{theorem}

Fix \(e=\min(\varepsilon,1/10)\) and suppose \(H\ge R^e\).
For large \(R\), the side of least radical contains a prime power of logarithm \(>2eH\).
The ratio of the other sides is then close to a sign at that prime.
Adjoining auxiliary primes gives a lattice of principal units with large first minimum.
Two determinant estimates force a nonzero polynomial to vanish on a slice of an integer simplex.

Write \(\lambda\) for the slope of the slice.
The first minimum excludes small denominators for its Fourier frequencies.
Steuding and Technau's prime estimate locates a prime \(q\) with
\(q\lambda\) close to \((0,1/m,\ldots,(m-1)/m)\)
\cite[Theorem~3]{SteudingTechnau}.
In each residue class modulo \(q\), the slice thus contains consecutive colors of a smaller simplex.
Finite differences give interpolation on those colors.
Lucas's congruence lifts interpolation to the full slice, contradicting the vanishing polynomial.

The use of AI tools is described at the end of the paper.


\section{Parameters and arithmetic estimates}\label{sec:inputs}

Order the summands as $a\le b$. Write $H=\log c$ and $3x=\log R$.
All logarithms are natural. We fix $e=\min(\varepsilon,1/10)$ and
the parameters of Table~\ref{tab:parameters}. Every threshold for $x$
is uniform in the triple.
\fixedparameters

Every triple satisfies $H\le C_0R^{b_0'}$ for absolute $C_0>0$,
$b_0'\ge0$ \cite[Proposition~2.1]{Z29}. Taking exponent $2/9+1/18$
in \cite[Theorem~1.1]{Z29} gives $\log H\le 5x/6+O(1)\le x$
for large $x$.

Use the conventions of \cite[Section~2]{Z29}. A reduced fraction
$A/B$ with $B>0$ has height $h(A/B)=\log\max(|A|,B)$.
We normalize $v_p(p)=1$. Fix $\vartheta>0$. For $x\ge x_\vartheta$,
each signed quotient $z_0\ne\pm1$ of two sides and each prime $q\le R$ satisfy
\begin{equation}\label{eq:local}
 -\log|1-z_0|\le \exp(\vartheta x),\qquad
 v_q(1-z_0)\log q\le q\exp(\vartheta x).
\end{equation}
These bounds are \cite[Section~3, Fact~F2]{Z29}, from
\cite[Fact~F3, Lemma~2.1 and (2.8)]{ZhangThird}.

\begin{fact}\label{fact:bl}
Let $p$ be odd and let $\alpha_1,\alpha_2$ be multiplicatively
independent rational principal units at $p$. Put $\ell=\log p$ and
$T_i=\max(h(\alpha_i),\ell)$. For positive integers $b_1,b_2$, put
$b'=b_1/T_2+b_2/T_1$. Then
\begin{equation}\label{eq:bl}
 v_p(\alpha_1^{b_1}-\alpha_2^{b_2})\le
 \frac{24p T_1T_2}{(p-1)\ell^4}
 \max\{\log(b'\ell)+0.4,\,10\ell,\,10\}^2 .
\end{equation}
\end{fact}
This is the rational principal-unit case of
\cite[Corollary~2]{BL96}, used in \cite[proof of Proposition~6.1]{Z29}.

\begin{fact}\label{fact:primes}
Write $\vartheta(z)=\sum_{p\le z}\log p$. For $M,X\ge10^9$,
\[
 \vartheta(M)-\vartheta(M/2)>M/3,\qquad
 \vartheta(X)-\vartheta(10X/11)>X/12 .
\]
\end{fact}
Both inequalities follow by applying
$|\vartheta(z)-z|<0.2z/(\log z)^2$ at the two endpoints
\cite[Theorem~5.2]{Dusart}.

\section{A prime of large contact}\label{sec:contact}

Put $3x=\log R$, order $a\le b$, and suppose $H\ge\exp(3ex)$ with $x$ sufficiently large.
Use the least-radical side $S$ and signed ratio $z$ of \cite[Section~2]{Z29}.
The three side radicals have product $R$, so $\rad(S)\le R^{1/3}=\exp x$.
Its rational logarithmic height is $H_z=h(z)\in[H-\log2,H]$, and
$v_p(z-1)=v_p(S)$ for $p\mid S$.
Apply \eqref{eq:local} with $\vartheta=3g$ as in
\cite[Propositions~3.1 and~4.1]{Z29}. At infinity $b/c$ gives
$\log(c/a)\le EH$, where $E=\exp(-3(e-g)x)$.
At $p\mid S$ with $p\le HR^{-2g}$, the ratio $z$ gives
$v_p(S)\log p\le H\exp(-3gx)$. There are at most $x/\log2$ such primes.
Since $\log S\ge(1-E)H$, the remaining primes have total mass
$\sum\rho_p>0.9$, where $\rho_p=v_p(S)\log p/H$.
To adapt \cite[Proposition~5.1]{Z29}, put $\omega_p=\log p/(3x)$.
Each $\omega_p>e-2g$ and $\sum\omega_p\le1/3$. A maximizer of
$\rho_p/\omega_p$ satisfies
\begin{equation}\label{eq:contact-prime}
 \rho_p>2.7(e-2g)>2e,\qquad 2.94ex<\ell=\log p\le x .
\end{equation}

In particular $p>HR^{-2g}\ge\exp(3(e-2g)x)>2$ for large $x$.
On the support $\mathcal S$ of $z$, let $\Lambda$ be the kernel of the residue
map $\beta(v)=\prod_i i^{v_i}$ into $C_p=\F_p^\times/\{\pm1\}$.
Its injective character $\gamma(v)$ is the signed product congruent to $1\bmod p$
\cite[Section~6]{Z29}. Use the same maps on adjoining primes away from $p$.
The height norm $N$ is the larger of the positive and negative exponent sums
weighted by $\log i$ \cite[Lemma~7.2, (7.4)]{Z29}.
Since $z\equiv1\bmod p$, its vector belongs to $\Lambda$ and has character $z$.
Write it as $\nu_pv_0$ with $\nu_p>0$ and $v_0$ primitive in $\Lambda$.
The scalar $h=N(v_0)=H_z/\nu_p$ is the height of this primitive vector.
The primitive-line argument of \cite[proof of Proposition~6.1]{Z29} gives
$N(v)\ge\ell-\log2$ for $v\ne0$ and
$v_p(S)\ell\le H/\nu_p+\log(2\nu_p)$.
Hence $\nu_p\le H/(\ell-\log2)\le H$.
By \eqref{eq:contact-prime} and $\log(2H)<eH$, we have $\nu_p\le1/e<p$.
Lifting at $p$ yields $n=v_p(\gamma(v_0)-1)=v_p(S)$ and
\begin{equation}\label{eq:contact-height}
 h\ge(e/2)\exp(3ex),\qquad
 2e<\delta=n\ell/h\le1+(\log2)/h .
\end{equation}
The numerator of $\gamma(v_0)-1$ gives the upper bound. At infinity,
\[
 |\log|\gamma(v_0)||/h\le\log(c/a)/H_z\le2E .
\]

\section{Auxiliary primes and the first minimum}\label{sec:aux}

Choose $i_*$ maximizing $M_*=|(v_0)_{i_*}|$.
The weighted absolute length of $v_0$ satisfies
$\sum_i|(v_0)_i|\log i\ge2(1-E)h$, so $M_*\ge h/(2x)$.
Project along $v_0$ onto $y_{i_*}=0$ by $\pi$.
Put $\|y\|_\perp=\sum_{i\ne i_*}|y_i|\log i$ and $u=x/\log x$.

\begin{lemma}\label{lem:separation}
Every nonzero point of $\pi\Lambda$ has transverse norm at least $2u$.
\end{lemma}
\begin{proof}
Suppose $0<\|\pi v\|_\perp<2u$ for $v\in\Lambda$.
Normalize area so that $\Z v_0+\Z v$ has covolume one, and let $A$ be the
area of the unit $N$-ball in its span. Take a basis $w_1,w_2$ attaining its successive minima
$L_i$ as in \cite[proof of Proposition~6.1]{Z29}, and write $v_0=\sum c_iw_i$.
The diamonds on $v_0/h$ and either $\pi v/N(\pi v)$ or $w_{3-i}/L_{3-i}$
lie in the ball by symmetry and convexity. Their areas and Minkowski's bound give
\[
 \max\left\{\frac{2}{hN(\pi v)},\frac{2|c_i|}{hL_{3-i}}\right\}
 \le A\le\frac4{L_1L_2},\qquad L_1L_2<4hu,\quad |c_i|L_i\le2h .
\]
Both $L_i\ge\ell-\log2$. A zero $c_i$ would give $L_{3-i}=h$ by primitivity
and hence $L_i<4u<\ell-\log2$.
Apply \eqref{eq:bl} to the two independent basis units, inverted as needed.
Here $T_i\le2L_i$, $b'\le4h/(L_1L_2)$ and $\log h\le x\ll_e\ell$.
Thus $n\ell\ll_e L_1L_2/\ell\ll_e h/\log x$, contradicting \eqref{eq:contact-height}.
\end{proof}

Set $M=\exp(e_0x)\max(\sqrt p,p/h)$, $t=\log M$ and $\tau=t/\ell$.
Since $\log h\ge3ex-\log(2/e)$ and $2.94ex<\ell\le x$, we get
\begin{equation}\label{eq:tau-red}
 \tau\le\max\{\tfrac12+e/29.4,\ 1-3e+e_0+\log(2/e)/\ell\}<1-2.8e .
\end{equation}
In particular $e_0x\le t\le x$ and $M<p/2$.
Fact~\ref{fact:primes} leaves $\gg M/t$ primes in $(M/2,M]$ after removing divisors of $2R$.
They are prime to $p$. Put $G=\beta(\Z^{\mathcal S})$, $s=|\mathcal S|-1\ge1$
(the complementary sides are coprime and exceed one). Take one torus
$\R^s/\pi\Lambda$ per coset of $G$ in $C_p$, as in
\cite[proof of Lemma~6.2, overlap count]{Z29}.
Fix a representative per coset; for each $q$ choose $a_q\in\Z^{\mathcal S}$
so $q\beta(a_q)$ equals its coset representative, and mark $\pi a_q$ in that torus.
The covolume is $|G|/M_*$ by expanding a lattice basis starting with $v_0$
along coordinate $i_*$ after projection of the other columns. Thus the total
volume $\mu$ and the volume $W_0$ of an embedded ball of radius $u/2$ satisfy
\[
 \mu=\frac{p-1}{2M_*}\le\frac{xp}{h},\qquad
 W_0=\frac{u^s}{s!\prod_{i\ne i_*}\log i}\ge(u/(3x))^s=\exp(-o(x)).
\]
Here Lemma~\ref{lem:separation} gives the embedding, and the volume bound uses
$\sum_i\log i\le3x$, $s=O(x/\log x)$ as in \cite[proof of Lemma~7.2, transverse volume]{Z29}.
Averaging gives one ball containing $n_c\gg MW_0/(t\mu)\ge\exp(e_0x-o(x))$ primes.
Lift the marked points in this ball to a single ball in $\R^s$.
Their pairwise differences are $\pi o(q,q')$ for choices
\[
 o(q,q')\in a_q-a_{q'}+\Lambda,\qquad
 \|\pi o(q,q')\|_\perp<u,\qquad \beta(o(q,q'))q/q'=1\quad\text{in } C_p .
\]

\begin{proposition}\label{prop:minimum}
Some distinct $q_0,\ldots,q_m$ in this ball give first minimum
$J>\exp(\kappa x)$ for the lattice $\Gamma_d$ generated by $v_0$ and
$u_j=o(q_j,q_0)+\mathbf e_{q_j}-\mathbf e_{q_0}$, $1\le j\le m$.
\end{proposition}
\begin{proof}
Put $Y=\exp(\kappa x)<h$. In a nonzero vector $a'v_0+\sum b_j u_j$
of norm at most $Y$, the $d$ fresh coefficients sum to zero.
Their absolute values are at most $Y$. At least one is nonzero.
Fix them all, and let $i\ge1$ be the last nonzero index.
Fix all primes except $q_i$, the old $i_*$-coordinate, and $a'\bmod |b_i|$.

For two completions, the difference of the old parts divided by $b_i$
is integral and vanishes at $i_*$. Its transverse norm is $<2u$
by cancellation of the fixed lifts. The divisor count in
\cite[proof of Lemma~6.2, small-denominator count]{Z29}
allows at most $(1+2u)\exp(2u)$ differences from a fixed completion.
Relative to a fixed completion $\bar q_i$, the divided difference $w'$ equals
$o(q_i,q_0)-o(\bar q_i,q_0)$ plus an integral multiple of $v_0$.
Thus $w'+\mathbf e_{q_i}-\mathbf e_{\bar q_i}$ has trivial signed residue.
Each $w'$ fixes the signed class of $q_i$, with at most one prime in $(M/2,M]$ since $M<p/2$.
The zero sum allows $O_e(Y^{d-1})$ coefficient tuples.
The old coordinate and axial residue each have $O(Y)$ choices.
Among $(1-o(1))n_c^d$ distinct tuples the bad proportion is therefore
$O_e(Y^{d+1}(1+2u)\exp(2u)/n_c)\le\exp(-2\kappa x+o(x))$,
since $(d+3)\kappa=e_0$.
\end{proof}

\section{A box and a vanishing polynomial}\label{sec:box}

Fix the tuple of Proposition~\ref{prop:minimum}. Reduce its generators modulo $v_0$:
$u_j=\lambda_jv_0+\pi_j+\mathbf e_{q_j}-\mathbf e_{q_0}$,
where $|\lambda_j|\le1/2$ and $\|\pi_j\|_\perp<u$.
For $P=a'v_0+\sum y_j u_j$, put $Z_0=h(a'+\lambda\cdot y)$.
Set $a_1=(1-e/5)n\ell/[d(1+1/k)]$ and $2a_0=da_1/k$.
Define $\mathcal V$ by $|Z_0|\le a_0$, $ty_j\ge-a_1$ and $t\sum y_j\le a_1$.
Put $b_0=\lfloor a_1/t\rfloor$, $L=db_0$, $\theta=-b_0\sum\lambda_j$ and $W=2a_0/h$.
Write $T_m(U)=\{z'\in\Z_{\ge0}^m:|z'|_1\le U\}$.
Then $z'=y+b_0\mathbf1$ runs through
\[
 \mathcal B(L),\qquad
 \mathcal B(U)=\{z'\in T_m(U):\|\theta+\lambda\cdot z'\|_{\R/\Z}\le W/2\}.
\]
By \eqref{eq:contact-height}, $W=(1-e/5)\delta/(k+1)\in[w,1)$.
The circle condition gives an integer $a'$, unique since $W<1$.

Put $D=\lceil(1+c_0)n\rceil-1$, $X=\lceil1.1D/(k+1)\rceil$ and $L_*=D+X$.
Then $n\le D<(1+c_0)n$ and \eqref{eq:tau-red} gives
\begin{equation}\label{eq:box-margin}
 L/D\ge\frac{1-e/5}{(1+1/k)(1+c_0)\tau}-d/D>1+1.9e .
\end{equation}
Here $1/k\le e/2$ and $d/D=o(1)$ by $n\ge e^2\exp(3ex)/x$.
Hence $L_*\le L$.

Choose a shortest vector $v_1$. For $Q=\lfloor(h/t)^{1/d}\rfloor$, two phases
$\lambda\cdot y$, $y\in\{0,\ldots,Q\}^m$, have circle distance at most $(Q+1)^{-m}$.
Their difference with the nearest axial correction is nonzero.
Each transverse generator costs at most $t+u\le2t$. Thus
\begin{equation}\label{eq:short-upper}
 J\le h(Q+1)^{-m}+2mtQ\le(1+2m)h\sqrt{t/h}.
\end{equation}

\begin{proposition}\label{prop:vanish}
A nonzero rational polynomial of total degree at most $D$ vanishes on $\mathcal B(L)$.
\end{proposition}
\begin{proof}
Index columns by the coordinates $(a',y)$ of points in $\mathcal V\cup(\mathcal V+v_1)$
in the basis $v_0,u_1,\ldots,u_m$.
The first $N_0=\binom{n+m-1}{m}$ rows evaluate
$P_\beta(y)=\prod_i\binom{y_i}{\beta_i}$ for $|\beta|<n$.
The next $N_1=\lfloor c_0^mN_0\rfloor$ rows evaluate $\gamma(P)P_\beta(y)$
for the first $N_1$ indices in increasing degree.
For $V_1=\lceil c_0n\rceil$, the second block lies in $|\beta|<V_1$,
since $V_1+i\ge c_0(n+i)$ for $0\le i<m$.
Its mean degree is at most the full simplex mean $m(V_1-1)/d<mc_0n/d$.
Write $N_*=N_0+N_1$.

Translate the columns to nonnegative coordinates. By
\cite[proof of Proposition~7.1, valuation argument]{Z29},
binomial translation is unitriangular on both downward-closed blocks.
The translated character acquires a common $p$-unit factor.
After removing it, expand $\gamma(v_0)^{a'}\prod_i\gamma(u_i)^{y_i}P_\beta(y)$
in the basis $\binom{a'}{\mu}P_\nu(y)$ with coefficients $c_{\mu,\nu}$. Then
\[
 v_p(c_{\mu,\nu})\ge n\mu+\sum_i(\nu_i-\beta_i)_+ .
\]
Indeed, the axial unit has depth $n$. At each transverse unit the term of
binomial index $j$ has valuation at least $j$, and multiplying by
$\binom{y_i}{\beta_i}$ gives integral binomial coefficients supported at
$\nu_i\le j+\beta_i$. Only finitely many indices contribute on the columns.
Subtract terms with $\mu=0$, $|\nu|<n$ using the first block.
Each remaining entry has valuation at least $n-|\beta|$.
Summing over rows, any nonzero minor $\Delta$ of order $N_*$ satisfies
\begin{equation}\label{eq:det-lower}
 v_p(\Delta)\ell>N_1n\ell(1-mc_0/d).
\end{equation}

For the global estimate translate only the polynomial coordinates.
In $\mathcal V$, $|y_i|\le ma_1/t=O_e(n)$.
Translation by $v_1$ changes each $y_i$ by at most $J/\log(M/2)=O_e(n)$
by \eqref{eq:short-upper} and \eqref{eq:contact-height}.
Shifting each coordinate by its minimum therefore bounds the binomial entries by $\exp(O_e(n))$.
Every determinant term contains $N_1$ character factors. At each place take
the maximum over the union of both boxes.
The product formula and integrality of the binomial entries at finite places give
\[
 v_p(\Delta)\ell
 \le N_1\sum_{v\ne p}\max_{P\in\mathcal V\cup(\mathcal V+v_1)}
      \log|\gamma(P)|_v+O_e(N_*(n+\log N_*)).
\]
On $\mathcal V$ the old finite places contribute $\le2a_0+O_e(a_1u/t)$
since $v_0$ has weighted absolute length $\le2h$.
Each fresh place contributes $\le a_1$ by the defining inequalities.
Infinity contributes $\le2a_0E+O_e(a_1(u+\log2)/t)$ since $1/2<q_j/q_0<2$.
The translate adds at most $N(v_1)=J$ to the sum of maxima.
Now $2a_0+da_1=(1-e/5)n\ell$, $N_*/N_1=O_e(1)$ and $\log N_*=O_e(\log n)$.
Dividing by $N_1n\ell$ gives the upper bound
$1-e/5+O_e(E+(u+1)/t+\sqrt{t/h}+1/\ell)=1-e/5+o(1)$.
This contradicts \eqref{eq:det-lower}.

All minors of order $N_*$ vanish. With fewer than $N_*$ columns the rank is
already smaller than $N_*$. Thus a row dependence gives $P_0(y)+\gamma(P)P_1(y)=0$ on both boxes,
with degrees $\le n-1,V_1-1$. If either polynomial is zero, the other suffices.
Otherwise let $s$ be the displacement of $v_1$ in the $y$ coordinates.
Translate-elimination \cite[proof of Proposition~7.1, rank argument]{Z29}
gives $P_0(y+s)P_1(y)-\gamma(v_1)P_0(y)P_1(y+s)=0$ on $\mathcal V$.
Its top form is the product of the two top forms times $1-\gamma(v_1)\ne0$.
Its degree is at most $n+V_1-2\le D$. Translate to $z'$.
\end{proof}

\section{A prime for the selected slope}\label{sec:hit}

Relabel the transverse coordinates in their original order, from $0$ to $m-1$.
Put $\lambda_i^*=i/m$ and write $\|s\|$ for the distance to the nearest integer.

\begin{lemma}\label{lem:hit}
Some prime $q\in I=(10X/11,X]$ satisfies $\max_i\|q\lambda_i-\lambda_i^*\|<\eta$,
$D/(k+1)<q\le L_*-D$ and $q>k+1$.
\end{lemma}

\begin{proof}
If $v\in\Z^m$ and $0<|v|_1\le J/(3t)$, rounding $-\lambda\cdot v$
to the nearest axial integer gives a nonzero vector of norm at most
$h\|\lambda\cdot v\|+(t+u)|v|_1$. Thus $\|\lambda\cdot v\|\ge J/(2h)$
by $u<t/2$ and minimality of $J$.
Write $T=\log X$. From $X\asymp_e n$, \eqref{eq:contact-height},
$\ell\le x$ and the upper bound on $\log H$, we have $T\asymp_e x$ and
$h/X\ll_e x$. Hence $J>\exp(\kappa x)$ and $t\le x$ give
$mKT^{20}<J/(3t)$ and $T^{20}/X<J/(2h)$.

For $0<|b|_\infty\le K$, $b\in\Z^m$, apply Dirichlet approximation to
$\beta=b\cdot\lambda$ with $Q=\lfloor X/T^{20}\rfloor$. It gives
coprime $A,r$, $1\le r\le Q$, with $|r\beta-A|\le1/(Q+1)<T^{20}/X$.
A denominator $r\le T^{20}$ contradicts the lower bound at $v=rb$.
Thus $T^{20}<r\le X/T^{20}$ and $|\beta-A/r|<r^{-2}$.

Write $\Lambda_{\rm vM}$ for the von Mangoldt function and $\chi$ for the
periodic indicator of $(-1/4,1/4]$. At $N=\lfloor X\rfloor,\lfloor10X/11\rfloor$,
\cite[Theorem~3, (15)]{SteudingTechnau} with exponent $1/100$ gives uniformly in $s$
\[
 \left|\sum_{j\le N}\Lambda_{\rm vM}(j)(\chi(j\beta-s)-1/2)\right|
 \ll T^8(X/\sqrt r+X^{3/4}+\sqrt{Xr}+X^{33/40}),
\]
since $(4Nr)^{3/400}N^{81/100}\ll X^{33/40}$.
The bounds on $r$ give $O(X/T^2)$. Integrating against
$\exp(2\pi i s)$ over one period replaces the centered arc by
$\exp(2\pi i j\beta)/\pi$. Subtraction bounds the exponential sum on $I$ by $O(X/T^2)$.

Set $F(s)=\prod_i\eta^{-1}(1-\|s_i-\lambda_i^*\|/\eta)_+$.
Each factor has mean one and Fourier coefficients of modulus
$(\sin(\pi a\eta)/(\pi a\eta))^2$ for $a\ne0$. Their absolute sum is
$1/\eta$ by evaluating the unshifted series at zero. The tail beyond $K$
is at most $2/(\pi^2K\eta^2)$ by summing $(\pi a\eta)^{-2}$.
Thus the product has Fourier mass $\eta^{-m}$, with at most
$2m/(\pi^2K\eta^{m+1})\le1/4$ outside the cube. Put
$S_I=\sum_{j\in I}\Lambda_{\rm vM}(j)>X/12$ by Fact~\ref{fact:primes}.
Higher prime powers have at most $T/\log 2$ exponents and $\sqrt X$ bases per exponent, so
\[
 \sum_{q\in I,\ q\ {\rm prime}}(\log q)F(q\lambda)
 \ge \tfrac34 S_I-O\bigl(\eta^{-m}(X/T^2+\sqrt X\,T^2)\bigr)>0 .
\]
Any contributing $q$ has the required distances. Also $10X/11\ge D/(k+1)$
and $X=L_*-D\asymp_e D\to\infty$, giving the remaining inequalities.
\end{proof}

\section{Interpolation on a simplex}\label{sec:interpolation}

Give each point \(y\in T_m(U)\) the color
\(\sigma(y)=\sum_{i=0}^{m-1}i y_i\bmod m\), so
\(\lambda^*\cdot y\equiv \sigma(y)/m\pmod 1\).

\begin{lemma}\label{lem:color}
Let \(q\) be prime and \(0\le j<\min(m,q)\).
Choose any \(j+1\) consecutive colors modulo \(m\).
Their points in \(T_m(j+1)\) determine every polynomial over \(\F_q\)
of total degree at most \(j\).
\end{lemma}

\begin{proof}
Read colors modulo \(m\). Suppose \(P\) vanishes on the colors
\(a,\ldots,a+j\), with \(0\le a<m\). We induct on
\(j\). For \(j=0\), evaluation at the unit vector \(e_a\) detects
the constant. For \(j\ge1\) and \(0\le i<m-1\), consider
\[
  P(y+e_{i+1})-P(y+e_i).
\]
This difference has degree at most \(j-1\). It vanishes on the \(j\)
colors \(a-i,\ldots,a-i+j-1\) in \(T_m(j)\). Both translates then have
colors in the original interval. By induction the difference
is identically zero. Thus \(P\) is invariant under translation by each
\(e_{i+1}-e_i\), and hence by every \(e_i-e_0\).

Use coordinates \(s=\sum_i y_i,y_1,\ldots,y_{m-1}\).
Translation by \(e_i-e_0\) now increments only \(y_i\).
Suppose the leading term in \(y_i\) is \(A y_i^b\), with \(b>0\).
Its forward difference has leading term \(bA y_i^{b-1}\).
Since \(b\le j<q\), this contradicts invariance.
Consequently \(P=f(s)\). For \(1\le j'\le j+1\) the point
\((j'-1)e_0+e_a\) has color \(a\) and coordinate sum \(j'\).
These \(j+1\) distinct elements of \(\F_q\) are roots of \(f\),
so \(f=0\).
\end{proof}

\begin{lemma}\label{lem:block}
Fix the prime $q$ of Lemma~\ref{lem:hit}. For every
$r\in\{0,\ldots,q-1\}^m$ with $|r|_1\le D$, put
$j=\lfloor(D-|r|_1)/q\rfloor$. Evaluation on
$\{y\in\Z_{\ge0}^m:qy+r\in\mathcal B(L_*)\}$ is injective
on the polynomials of total degree at most $j$ over $\F_q$.
\end{lemma}

\begin{proof}
By Lemma~\ref{lem:hit}, $q>D/(k+1)$ and $q>k+1$, so $0\le j\le k<q$.
Choose a real lift of $\theta_r=\theta+\lambda\cdot r$. Put
$a=\lceil-m(\theta_r+w/4)\rceil$. For $0\le i\le j$ we have
\[
 -w/4\le\theta_r+(a+i)/m<-w/4+(j+1)/m\le w/4 .
\]
Here $(j+1)/m\le w/2$ by the choice of $m$.
Consider $y\in T_m(j+1)$ of color $a+i\bmod m$ for some
$0\le i\le j$. Then $\|\theta_r+\lambda^*\cdot y\|\le w/4$.
Uniformly in $r$, we have
$\|(q\lambda-\lambda^*)\cdot y\|\le |y|_1\eta\le(k+1)\eta=w/4$. Thus
$\|\theta+\lambda\cdot(qy+r)\|\le w/2\le W/2$.
Also $|qy+r|_1\le q(j+1)+|r|_1\le D+q\le L_*$. The block contains
all points of these colors in $T_m(j+1)$. Since $j\le k<m$
and $j<q$, Lemma~\ref{lem:color} gives injectivity.
\end{proof}

\begin{proposition}\label{prop:lucas}
Evaluation on \(\mathcal B(L_*)\) is injective on the space of rational
polynomials in \(m\) variables of total degree at most \(D\).
\end{proposition}

\begin{proof}
For a multi-index \(a\), write \(|a|=\sum_i a_i\), \(a!=\prod_i a_i!\)
and \(B_a(Z)=\prod_i\binom{Z_i}{a_i}\).
The \(B_a\) with \(|a|\le D\) form a rational basis by their leading terms
\(Z^a/a!\). Their evaluations are integers. Suppose their columns satisfy
a relation with coefficients \(\xi_a\in\F_q\). For each residue
\(s\in\{0,\ldots,q-1\}^m\) with \(|s|\le D\), put
\(j_s=\lfloor(D-|s|)/q\rfloor\) and
\(F_s=\sum_{|v|\le j_s}\xi_{qv+s}B_v\).
Here \(j_s\le k<q\), so each \(v!\) is invertible.
At every node \(qy+r\) with \(|r|\le D\), Lucas's congruence~\cite{Lucas} gives
\(\sum_{s\le r} B_s(r)F_s(y)=0\), with coordinatewise order.
Indeed, compare coefficients of \(T^{qv_i+s_i}\) in
\((1+T)^{qy_i+r_i}\equiv(1+T^q)^{y_i}(1+T)^{r_i}\pmod q\).

Induct on \(|r|\). The terms with \(s\ne r\) vanish by induction,
and \(B_r(r)=1\). Thus \(F_r\) vanishes on the \(r\)-block.
Lemma~\ref{lem:block} gives \(F_r=0\) since \(\deg F_r\le j_r\).
The distinct leading terms \(Y^v/v!\) give \(\xi_{qv+r}=0\) for \(|v|\le j_r\).
Every column occurs, since the residue of \(a\) has sum at most \(|a|\).
Full column rank modulo \(q\) gives an integer maximal minor prime to \(q\),
hence full column rank over \(\Q\).
\end{proof}

\begin{proof}[Proof of Theorem~\ref{thm:main}]
Fix \(e=\min(\varepsilon,1/10)\). Let \(x_e\) exceed all preceding thresholds.
Suppose \(x\ge x_e\) and \(H\ge R^e\).
Proposition~\ref{prop:vanish} supplies a nonzero rational polynomial of
degree at most \(D\) vanishing on \(\mathcal B(L)\).
By \eqref{eq:box-margin}, \(L_*\le L\), so this polynomial vanishes on
\(\mathcal B(L_*)\). Proposition~\ref{prop:lucas} forces it to be zero.
Thus \(H<R^e\) for \(x\ge x_e\). For \(x<x_e\),
the cap in Section~\ref{sec:inputs} bounds \(H\) by a constant depending only on \(e\).
Since \(1\le R^e\le R^\varepsilon\), these bounds give the theorem.
\end{proof}
\section*{Use of AI tools}
The author devised the method of proof. Codex (OpenAI, GPT-6 Astra) and
Antigravity (Google, Gemini 4 Neon) ran blind reviews of the manuscript,
and Codex assisted with the Lean formalization. Claude Code (Anthropic,
Claude Fable 5) checked the proofs and the references. The author checked
every step of the proof and every reference, and takes full responsibility
for the content of the paper.

\section*{Statements and declarations}
{\raggedright
\textbf{Competing interests.} The author has no competing interests.\par
\textbf{Funding.} No funding was received for this work.\par
\textbf{Data availability.} A Lean~4 formalization of Theorem~\ref{thm:main}
is available at \url{https://github.com/ryzfsu/abc-log-bound-subexponential}
under the declaration \texttt{AbcSubexp.subexp}.\par}
\begin{thebibliography}{9}

\bibitem{BL96}
Y.~Bugeaud and M.~Laurent, \emph{Minoration effective de la distance
\(p\)-adique entre puissances de nombres alg\'ebriques}, J. Number Theory
\textbf{61} (1996), 311--342.

\bibitem{Dusart}
P.~Dusart, \emph{Estimates of some functions over primes without R.H.},
preprint, 2010, \url{https://arxiv.org/abs/1002.0442}.

\bibitem{Lucas}
\'E.~Lucas, \emph{Th\'eorie des fonctions num\'eriques simplement
p\'eriodiques}, Amer. J. Math. \textbf{1} (1878), 184--196, 197--240,
289--321.

\bibitem{PastenSep}
H.~Pasten and R.~Sep\'ulveda-Manzo, \emph{On the \(abc\) and the
\(abcd\) conjectures}, Bull. Braz. Math. Soc. (N.S.) \textbf{56}
(2025), Art.~54. Preprint version \url{https://arxiv.org/abs/2406.05083v1}.

\bibitem{SteudingTechnau}
J.~Steuding and M.~Technau, \emph{The least prime number in a Beatty
sequence}, J. Number Theory \textbf{169} (2016), 144--159.
Equation numbers follow \url{https://arxiv.org/abs/1512.08382v2}.

\bibitem{SY2001}
C.~L.~Stewart and K.~Yu, \emph{On the \(abc\) conjecture, II}, Duke Math. J.
\textbf{108} (2001), 169--181.

\bibitem{ZhangThird}
R.~Zhang, \emph{Improving the Stewart--Yu bound for all \(abc\) triples},
Zenodo, 2026. \url{https://doi.org/10.5281/zenodo.24202155}.

\bibitem{Z29}
R.~Zhang, \emph{An upper bound for log c in the abc conjecture with exponent 2/9}, Zenodo, 2026. \url{https://doi.org/10.5281/zenodo.24405683}.
\end{thebibliography}

\end{document}
